Recall that $a\in S(\mathfrak{a})$ exactly when $sa\in\mathfrak{a}$ for some $s\in S$.

\textbf{(i)} One inclusion is immediate. If $sa\in\mathfrak{a}$ and $ta\in\mathfrak{b}$ with $s,t\in S$, then $sta\in\mathfrak{a}\cap\mathfrak{b}$, giving the other.

\textbf{(ii)} If $sa\in r(\mathfrak{a})$, then $s^na^n\in\mathfrak{a}$ for some $n$, so $a\in r(S(\mathfrak{a}))$. Conversely, if $sa^n\in\mathfrak{a}$, then $(sa)^n\in\mathfrak{a}$, so $a\in S(r(\mathfrak{a}))$.

\textbf{(iii)} Apply the membership criterion to $1$.

\textbf{(iv)} Membership in either side means $s_1s_2a\in\mathfrak{a}$ for some $s_1\in S_1$ and $s_2\in S_2$.

If $\mathfrak{a}=\bigcap_{i=1}^n\mathfrak{q}_i$, then $S(\mathfrak{a})=\bigcap_iS(\mathfrak{q}_i)$. A primary component saturates to itself when its radical is disjoint from $S$, and to $A$ otherwise. Thus at most $2^n$ ideals occur.
